\documentclass[letterpaper,10pt,conference]{ieeeconf}
\overrideIEEEmargins
\usepackage{amsmath,amsfonts,amssymb}
\usepackage{algorithmic}
\usepackage{algorithm}
\usepackage{array}
\usepackage{tabularx}
\usepackage{booktabs}
\usepackage{threeparttable}
\newcolumntype{Y}{>{\centering\arraybackslash}X}
\newcolumntype{L}{>{\raggedright\arraybackslash}X}
\usepackage{textcomp}
\usepackage{stfloats}
\usepackage{url}
\usepackage{verbatim}
\usepackage{graphicx}
\usepackage{cite}

\usepackage{bm}
\usepackage{pifont}
\newtheorem{lemma}{Lemma}
\newtheorem{proposition}{Proposition}
\newtheorem{theorem}{Theorem}
\newtheorem{corollary}{Corollary}
\newtheorem{remark}{Remark}
\usepackage{xcolor}
\newcommand{\rev}[1]{\textcolor{blue}{#1}}
\newcommand{\nd}[1]{\textcolor{red}{#1}}
\newcommand{\todo}[1]{\textcolor{red}{\textbf{[TODO: #1]}}}
\newcommand{\draftnote}[1]{\par\noindent\textcolor{red}{\textbf{$\blacktriangleright$ DRAFT: #1}}\par\medskip}
\newcommand{\figph}[2]{\fcolorbox{red}{gray!8}{\parbox[c][#1][c]{0.96\linewidth}{\centering
\color{red!65!black}\small\textbf{[figure placeholder]}\\[3pt] #2}}}
\graphicspath{{figs/}}
\hyphenation{op-tical net-works semi-conduc-tor IEEE-Xplore Koop-man}

% ---- method-name macros (freeze notation here) ----
\newcommand{\mbdtrue}{\textsc{MBD-true}}
\newcommand{\dkmbd}{\textsc{DK-MBD}}
\newcommand{\dksplit}{\textsc{DK-MBD-split}}
\newcommand{\bkmbd}{\textsc{BK-MBD}}

% \usepackage[hidelinks]{hyperref}
\usepackage{hyperref}
\usepackage{subcaption}
% (a)/(b) subfigure labels in Times New Roman (ptm = Times); caption text unchanged
\DeclareCaptionFont{timesnr}{\fontfamily{ptm}\selectfont}
\captionsetup[subfigure]{labelfont=timesnr}

\begin{document}
% \setlength{\abovecaptionskip}{0.25\abovecaptionskip} % 그림 캡션 위 간격
% \setlength{\belowcaptionskip}{0pt}

\title{Koopman-Accelerated Model-Based Diffusion\\for Real-Time Robot Control}

% Author intentionally omitted for blind submission.
% \author{Author Name(s)\thanks{Affiliation / funding.}}

\maketitle

\begin{abstract}
Model-Based Diffusion (MBD) provides derivative-free annealed sampling but requires thousands of rollouts per control step, making simulator-based MBD too slow for real-time control.
We replace the simulator with a learned bilinear deep Koopman surrogate suited to configuration-dependent input gains and use a quantile-fitted error tube for inexpensive candidate ranking.
Noise then controls both candidate reach and sampling dispersion, so wide stages explore while narrow stages settle.
On a physics-free, analytically integrated kinematic Franka Research 3 (FR3) testbed and a full rigid-body FR3 simulation in MuJoCo, the surrogate matches the corresponding true-dynamics oracle in all $70$ trials while meeting every $50$\,ms planning deadline, compared with $4.1$\,s per step for the simulator-in-the-loop oracle.
At equal sample budgets on a $1$\,cm stress test, fixed-noise schedules achieve $19$--$31$ of $35$ successes, whereas annealing achieves $34/35$ with roughly half the median settling time of fixed-wide sampling.
On non-convex wall reaching, the sampling planner threads the passage without violations, whereas convexified bilinear Koopman model predictive control trades passage completion against safety.
\nd{Finally, the planner, cost, and hyperparameters transfer unchanged to a physical FR3, tracking a continuously moving target in real time, with only the rollout model reidentified from on-robot data.}
\end{abstract}

\noindent\textbf{Keywords---}Koopman operator, sampling-based control,
model-based diffusion, data-driven control, model learning for control.

%==============================================================
\section{Introduction}\label{sec:intro}
%==============================================================
Sampling-based optimal control has become a standard approach for robotic systems with complex nonlinear dynamics.
Model Predictive Path Integral (MPPI) control~\cite{kappen2005, theodorou2010, williams2017mppi, williams2016aggressive} and Model-Based Diffusion (MBD)~\cite{pan2024model, xue2025full} evaluate hundreds of sampled control sequences by forward simulation and update the nominal sequence by a cost-weighted average.
MBD reinterprets this primitive as reverse diffusion, in which the weighted mean estimates the score of a Gaussian-smoothed target~\cite{song2019generative, song2020score} and a decreasing noise schedule yields coarse-to-fine search.
This derivative-free structure is general and parallel, but its rollout model is evaluated thousands of times per control step.
When those rollouts come from a physics simulator~\cite{pan2024model}, they dominate latency and usually confine MBD to offline optimization, since receding-horizon control repeats the entire population after every action.

The Koopman operator~\cite{koopman1931, mezic2005} provides a principled means of replacing expensive nonlinear rollouts.
Lifting through hand-designed~\cite{williams2015edmd, proctor2016dmdc} or learned~\cite{lusch2018deep} observables reduces nonlinear rollout to matrix--vector products whose cost depends on the lifted dimension rather than plant complexity.
Linear lifted predictors are common because they preserve convex Koopman MPC~\cite{korda2018mpc, strasser2026}, and a learned linear model has recently accelerated MPPI~\cite{mppidk2026}.
Their constant lifted input matrix is nevertheless mismatched to tracked outputs driven through configuration-dependent gains.

A prior representability analysis~\cite{kim2026bilinear} showed that configuration-dependent input gains, such as a manipulator Jacobian, produce irreducible dispersion error under a constant linear input map and identified the bilinear class as the smallest structural remedy.
In MPC, the bilinear term destroys convexity, so prior work freezes the lifted state over the horizon~\cite{bruder2021bilinear}.
A sampling planner imposes no convexity requirement and can instead roll out the bilinear model directly, retaining the state-dependent input channel within the same matrix--vector computational regime.

Unlike diffusion models used as \emph{learned} motion priors~\cite{janner2022planning, chi2025diffusion, carvalho2023motion}, our method does not learn a trajectory distribution, since diffusion serves only as an optimizer while learning is confined to the rollout dynamics.
Consequently, the same trained dynamics surrogate can support new goals and costs without retraining a policy or trajectory prior.

We ask whether this surrogate preserves MBD quality at real-time cost, and what annealing contributes when model error corrupts the evaluated landscape.
Section~\ref{subsec:anneal_theory} separates noise-dependent candidate \emph{reach} from the estimator's \emph{precision floor}, showing that wide stages explore across basins whereas narrow stages settle, a trade-off that no single fixed scale can resolve.

The contributions of this paper are as follows:
\begin{itemize}
\item \textbf{Koopman-accelerated MBD at the control rate.}
We combine bilinear Koopman rollout, a quantile-fitted tube, and MBD, and the surrogate matches the oracle in all $70$ trials and reduces the simulator-inclusive step from $4.1$\,s to $64$\,ms. The optimizer requires $17$\,ms median and $40$\,ms worst-case latency, with zero $50$\,ms deadline misses\nd{, and the planner transfers to a physical FR3 with only the rollout model reidentified on-robot}.
\item \textbf{Annealing under surrogate error.}
We identify and measure the two noise levers under one surrogate at equal budgets on a $1$\,cm stress test. Wide-fixed sampling stalls at its dispersion floor ($28$--$31$ of $35$), narrow-fixed sampling locks onto minima of the surrogate's own landscape centimeters from the target ($19/35$), and only annealing supplies both levers, reaching $34/35$ in nearly half the settling time.
\item \textbf{Beyond convex Koopman MPC.}
Convexified bilinear MPC matches MBD on quadratic reaching, but on non-convex wall reaching MBD threads the passage without violations while MPC trades passage completion against safety.
\end{itemize}

%==============================================================
\section{Preliminaries}\label{sec:prelim}
%==============================================================

\subsection{Problem Formulation}\label{subsec:problem}

Consider the discrete-time system
\begin{equation}
s_{t+1} = f(s_t, u_t), \qquad
s_t \in \mathcal{X} \subset \mathbb{R}^{n},\quad
u_t \in \mathcal{U} = [\underline{u}, \bar{u}]^{m},
\label{eq:sys}
\end{equation}
in the \emph{model-free, real-time} setting, where the dynamics $f$ are unknown or too costly to integrate at the sampling rate, so that a true-dynamics rollout is available only as an \emph{oracle} reference, not as a deployable controller.
The model is learned solely from interaction data $\mathcal{D} = \{(s_i, u_i, s_i^{+})\}$.

Each task defines observable features $b(s) \in \mathbb{R}^{d}$ used both as the Koopman training state and in the cost, and every fidelity level uses $b = [q, p_{\mathrm{tcp}}]$, the joint positions and tool-center-point position.
The controller operates in a receding horizon, so that at state $s_t$ it seeks a control sequence $U \in \mathbb{R}^{mT}$ minimizing
\begin{equation}
J(U; s_t) = \sum_{k=1}^{T} c\big(\hat{b}_k, u_{k-1}\big)
+ \phi\big(\hat{b}_T\big),
\label{eq:cost}
\end{equation}
where $\hat b_k$ is produced by the rollout model under study. The controller applies the first control $u_0$ to \eqref{eq:sys} and then replans.

\subsection{Model-Based Diffusion: A Review}\label{subsec:mbd}

MBD~\cite{pan2024model} treats \eqref{eq:cost} as sampling from the Boltzmann target
\begin{equation}
p_0(U) \;\propto\; \exp\!\big(-J(U)/\alpha\big)
\label{eq:target}
\end{equation}
and descends an annealing ladder of Gaussian-smoothed marginals
\begin{equation}
p_{\sigma}(U) = \big(p_0 * \varphi_{\sigma}\big)(U)
= \mathbb{E}_{\varepsilon \sim \mathcal{N}(0, \sigma^{2} I)}
\big[\, p_0(U - \varepsilon) \,\big],
\label{eq:smooth}
\end{equation}
with $\sigma_1 > \cdots > \sigma_S$.
By Tweedie's formula, the score of the smoothed target is a denoising mean shift,
\begin{equation}
\nabla_U \log p_{\sigma}(U)
= \frac{\mathbb{E}\big[\,U_0 \mid U\,\big] - U}{\sigma^{2}} .
\label{eq:tweedie}
\end{equation}
Because the smoothing kernel $\varphi_\sigma(U-U')$ is the density of the proposal $U' \sim \mathcal{N}(U,\sigma^2 I)$, the posterior mean in \eqref{eq:tweedie} is a ratio of expectations under this proposal. With $N$ candidates $U_k=U+\sigma\varepsilon_k$, $\varepsilon_k\sim\mathcal{N}(0,I)$, the self-normalized estimator is the cost-weighted mean
\begin{equation}
\widehat{\mathbb{E}}\big[U_0 \mid U\big] = \sum_{k=1}^{N} w_k U_k,
\quad
w_k = \frac{e^{-(J(U_k) - J_{\min})/\alpha}}
           {\sum_{j} e^{-(J(U_j) - J_{\min})/\alpha}} .
\label{eq:weights}
\end{equation}
Candidates are projected onto $\mathcal{U}^{T}$, using the same projection for all compared methods.
The update at stage $s$ is a preconditioned score step with optional Langevin noise, \emph{with the step size tied to the noise level}:
\begin{equation}
U \leftarrow U + \eta_s\, \widehat{\nabla_U \log p_{\sigma_s}}(U)
+ \sqrt{2 \eta_s}\, \xi,
\qquad
\eta_s = \eta\, \sigma_s^{2} .
\label{eq:update}
\end{equation}
With $\eta = 1$ and no Langevin noise, \eqref{eq:update} reduces exactly to the weighted-mean update $U \leftarrow \sum_k w_k U_k$.

%==============================================================
\section{Bilinear Koopman Model-Based Diffusion}\label{sec:method}
%==============================================================

\subsection{Deep Koopman Lifting and Rollout Models}\label{subsec:models}

All learned rollout models share a deep lift with an exact-decoder convention.
With $d = \dim b$,
\begin{equation}
z = \Psi_\theta(b) =
\begin{bmatrix} b \\ \psi_\theta\big(\tau(b)\big) \end{bmatrix}
\in \mathbb{R}^{r},
\qquad
C z = z_{1:d} = b,
\label{eq:lift}
\end{equation}
where $\psi_\theta$ is a small $\tanh$ network and $\tau$ is a fixed task input map ($\tau(b) = [\sin q, \cos q]$, $r = 20$ at every fidelity level).
Slicing the decoder avoids learning a fragile reconstruction and ties the first $d$ lifted coordinates to the physical output.
Two lifted dynamics classes are compared:
\begin{align}
\textbf{linear (DK):}\quad
& z^{+} = A z + B u ,
\label{eq:lin}\\[2pt]
\textbf{bilinear (BK):}\quad
& z^{+} = A z + B_0 u + \sum_{i=1}^{m} u_i B_i z ,
\label{eq:bil}
\end{align}
where \eqref{eq:bil} is compactly written as $z^{+} = M(u)\, z + B_0 u$ with the input-dependent transition matrix $M(u) = A + \sum_i u_i B_i$.
For control-affine systems the continuous-time Koopman generator carries a bilinear input structure~\cite{iacob2024}, and \eqref{eq:bil} is its finite-dimensional, discrete-time rollout approximation, restricting the input map's state dependence to be linear in the lifted state. Its input tangent map $\partial z^{+} / \partial u_i = B_0 e_i + B_i z$ depends on the state through $B_i z$, whereas the linear form's $\partial z^{+}/\partial u = B$ is constant.
Both classes roll out in the same matrix--vector regime, and at planning time the current output is lifted once, $z_0 = \Psi_\theta(b(s_t))$, after which every candidate is propagated entirely in the lifted space.

\subsection{The Bilinear Class as the Prescribed Rollout}\label{subsec:representability}

We select the rollout class using the prior representability criterion of~\cite{kim2026bilinear}.
When the tracked output moves as $\Delta y = \Phi(s)\, u\, \Delta t + O(\Delta t^{2})$ with a configuration-dependent gain $\Phi(s)$, here the manipulator Jacobian on the end-effector rows, every input-to-output Markov parameter $C A^{k} B$ of \eqref{eq:lin} is state-independent, and the best constant input map is the \emph{mean} of $\Phi$ along the visited states, with an irreducible one-step error equal to the gain's dispersion, a model-class property that no amount of training can remove.
The bilinear form restores exactly the discarded structure, because the columns $C B_i z$ of its input gain are linear in $z$, and the Jacobian's entries, multilinear in $\{\sin q_i, \cos q_i\}$ across joints, become expressible once the learned features supply the cross-products.
The criterion is a representability statement about the one-step input channel, not a closed-loop theorem. Section~\ref{subsec:exp_mechanism} traces how an optimizer amplifies the gap into closed-loop failure, and Sec.~\ref{subsec:exp_franka} tests the prescription under full rigid-body physics.

\subsection{Multi-Step Training}\label{subsec:training}

All learned models are trained on random-rollout snippets by the same $H$-step loss ($H = 15$),
\begin{equation}
\mathcal{L} = \frac{1}{H} \sum_{h=1}^{H}
\Big[ \big\| C \hat{z}_h - b_h \big\|^{2}
+ \gamma \big\| \hat{z}_h - \Psi_\theta(b_h)^{\mathrm{sg}} \big\|^{2}
\Big],
\label{eq:loss}
\end{equation}
with $\hat z_0 = \Psi_\theta(b_0)$, $\hat z_{h+1}$ given by \eqref{eq:lin} or \eqref{eq:bil}, $\gamma = 0.1$, and stop-gradient ($\mathrm{sg}$) on the latent target.
The bilinear input matrices are \emph{zero-initialized}, $B_i = 0$, so that training starts exactly at the linear model and adds bilinear structure only where the data demands it, which stabilizes the $m$ per-channel matrices on the $7$-DOF manipulator.
Each dataset is generated once and shared across all learned models for a given training seed, so the compared classes differ only in their lifted dynamics, \eqref{eq:lin} versus \eqref{eq:bil}.
One-step regression in place of \eqref{eq:loss} yields rollouts that diverge before the planning horizon.

\subsection{Model-Error Tube}\label{subsec:tube}

To make the planner aware of where the surrogate can be trusted, the one-step lifted residuals on the training set, $r_k = \Psi_\theta(b_{k+1}) - \big(M(u_k)\Psi_\theta(b_k) + B_0 u_k\big)$, are covered by a proportional bound $\|r_k\| \le c_x \|z_k\| + c_u \|u_k\|$ fitted at quantile coverage, where a least-squares fit of $(\hat c_x, \hat c_u)$ to the residual norms is rescaled by the empirical $(1-\delta)$-quantile of the coverage ratios,
\begin{equation}
\begin{aligned}
s_\delta &= Q_{1-\delta}\!\left[
\frac{\|r_k\|}{\hat c_x \|z_k\| + \hat c_u \|u_k\|}
\right], \\
(c_x, c_u) &= \max(s_\delta, 1) \,(\hat c_x, \hat c_u),
\end{aligned}
\label{eq:tubefit}
\end{equation}
with $\delta = 10^{-3}$.
Quantile coverage, rather than a worst-case fit over every residual, keeps the constants proportional to the typical model error, and Sec.~\ref{subsec:exp_cost} quantifies how a worst-case fit lets a few rest-to-motion transients dominate the planning cost.
During planning a scalar tube is propagated along each candidate with the precomputed norm bound $\bar m(u) = \|A\|_2 + \sum_i |u_i| \|B_i\|_2 \ge \|M(u)\|_2$:
\begin{equation}
e_{k+1} = \big(\bar m(u_k) + c_x\big) e_k
+ c_x \|\hat z_k\| + c_u \|u_k\|,
\qquad e_0 = 0 .
\label{eq:tube}
\end{equation}
Wherever the fitted bound holds, the recursion maintains $e_k \ge \|\delta_k\|$ for the true-vs-predicted lifted error $\delta_k = \Psi_\theta(b_k) - \hat z_k$, by induction on the exact error dynamics $\delta_{k+1} = M(u_k)\delta_k + r_k$ of the bilinear rollout.
The planner then minimizes the tube-augmented cost
\begin{equation}
\tilde J(U) = J(U) + \beta_e \sum_{k} e_k ,
\label{eq:tubecost}
\end{equation}
where $\beta_e$ scales with the task cost magnitude, and $\beta_e = 0$ recovers plain \bkmbd{}.

\subsection{Annealing under Surrogate Error}\label{subsec:anneal_theory}

With a learned rollout the sampler never sees the true objective, and every candidate is instead scored by
\begin{equation}
\hat J(U) \;=\; J(U) + \Delta J(U),
\label{eq:jhat}
\end{equation}
where $\Delta J$ collects the effect of model error and the tube term on the evaluated cost.
Ideally, a stage with noise level $\sigma$ descends a smoothed \emph{Boltzmann} landscape, since by \eqref{eq:tweedie} it ascends $\log p_\sigma$ with $p_\sigma = e^{-\hat J/\alpha} * \varphi_\sigma$, or equivalently descends the free energy
\begin{equation}
F_\sigma(U) \;=\; -\alpha \log\, \mathbb{E}_{\varepsilon \sim \mathcal{N}(0, \sigma^{2} I)}\,
e^{-\hat J(U + \varepsilon)/\alpha} .
\label{eq:freeenergy}
\end{equation}
Equation \eqref{eq:freeenergy} is the unconstrained idealization, whereas the deployed planner projects candidates onto $\mathcal{U}$, under which the Gaussian identities of Sec.~\ref{subsec:mbd} hold only approximately (at a cold-start center, $21\%$ of proposal coordinates clip at $\sigma = 1.2$ under the $1.5$\,rad/s limits), so what Sec.~\ref{subsec:exp_anneal} measures is accordingly the \emph{clipped-proposal} landscape the planner actually evaluates.
Rather than assert what smoothing does to this landscape's topology, we locate the schedule's leverage in two elementary properties of the update itself, and measure the rest.

\begin{proposition}[Directional reach of a stage]\label{prop:reach}
Let $U_k = U + \sigma \varepsilon_k$ with $\varepsilon_k \overset{\mathrm{iid}}{\sim} \mathcal{N}(0, I_{mT})$, $k = 1, \dots, N$, and let $U^{+} = \sum_k w_k U_k$ with any weights $w_k \ge 0$, $\sum_k w_k = 1$.
For every unit direction $v$ and $\rho \in (0, 1)$, with probability at least $1 - \rho$,
\begin{equation}
\big\langle U^{+} - U,\, v \big\rangle
\;\le\; \max_k\, \big\langle \sigma \varepsilon_k,\, v \big\rangle
\;\le\; \sigma \sqrt{2 \log (N/\rho)} .
\label{eq:reach}
\end{equation}
\end{proposition}
\begin{proof}
A convex combination is bounded by its largest element. Because $\langle \varepsilon_k,v\rangle$ are i.i.d.\ standard normal, a Gaussian tail bound followed by a union bound gives $\Pr[\max_k \langle \varepsilon_k,v\rangle > \sqrt{2\log(N/\rho)}] \le N e^{-\log(N/\rho)}=\rho$.
\end{proof}

The proposition bounds the pre-projection candidate cloud, whereas the deployed box projection can redirect displacements near the bounds and is treated empirically.
At the deployed $N = 800$ and $\rho = 10^{-3}$, the envelope is approximately $5.2\sigma$. For any fixed direction, a single update exceeds this envelope with probability at most $\rho$. Whether a basin within the envelope is reached depends on its geometry, which Secs.~\ref{subsec:exp_window} and~\ref{subsec:exp_anneal} examine empirically.
Nor can replanning substitute for reach across a cost barrier, because intermediate nominals are penalized.
Reach is bought with noise, and noise has a price at the other end.

\begin{remark}[Precision floor]\label{rem:floor}
The weighted-mean update at noise level $\sigma$ moves the nominal by $\sigma \sum_k w_k \varepsilon_k$, and under a locally frozen-weight approximation (the weights are computed from the same $\varepsilon_k$, so this is a scaling heuristic rather than an exact variance), it disperses per axis on the order of $\sigma / \sqrt{N_{\mathrm{eff}}}$, $N_{\mathrm{eff}} = 1 / \sum_k w_k^{2}$.
In receding horizon the executed control inherits this dispersion at every step, so a sampler at fixed $\sigma$ resists settling below it, and only annealing, which lowers the final noise level, reduces it at a fixed sample budget.
Section~\ref{subsec:exp_anneal} measures the dispersion directly by repeated planning from a fixed state.
\end{remark}

A fixed noise level must therefore trade reach against precision, whereas the anneal is the schedule that opens with the reach and closes at the floor, and Sec.~\ref{subsec:exp_anneal} validates each lever by its own failure signature.

In summary, one control step of \bkmbd{} lifts the measured output once by \eqref{eq:lift}. Then, for each annealing stage $s = 1, \dots, S$, it samples $N$ candidate sequences, propagates them in the lifted space by \eqref{eq:bil} while accumulating \eqref{eq:tube}, weights them by \eqref{eq:weights} applied to \eqref{eq:tubecost}, and updates the nominal sequence by \eqref{eq:update}. Finally, it executes the first control, warm-starts the next step by a one-step shift, and replans.

%==============================================================
\section{Experiments}\label{sec:exp}
%==============================================================
We evaluate end-effector reaching with one robot, the Franka Research 3 (FR3), at increasing fidelity levels.
The \emph{kinematic} FR3 runs with no physics engine, so its true dynamics are the analytic velocity integrator $q_{t+1} = \mathrm{clip}(q_t + u_t\,\Delta t)$ over the real FR3 joint axes, link transforms, and joint ranges, isolating model class and optimizer precision from every physics effect, and it carries the optimizer study.
The \emph{MuJoCo} FR3 adds full rigid-body physics and a joint-velocity servo that realizes the first-order input interface only up to lag and load\nd{, and the \emph{physical} FR3 tests real-time deployment (Sec.~\ref{subsec:exp_hw})}.
The experiments test the claims of Sec.~\ref{sec:intro} in sequence. Section~\ref{subsec:exp_franka} establishes oracle parity across rollout classes and optimizer families, Sec.~\ref{subsec:exp_window} examines the limit of a convexified alternative on wall reaching, Sec.~\ref{subsec:exp_mechanism} traces the amplification of a modest open-loop error into closed-loop failure, Sec.~\ref{subsec:exp_anneal} studies annealing at equal sample budgets and by direct spectrum measurement, and Sec.~\ref{subsec:exp_cost} reports rollout cost, latency, and tube coverage.



\begin{table}[t]
\centering
\begin{threeparttable}
\caption{Hyperparameters.}
\label{tab:hyper}
\footnotesize
\begin{tabularx}{\linewidth}{l L}
\toprule
Encoder $\psi_\theta$ & $\tanh$, $2$ layers $\times$ $96$ units; lift $r = 20$ \\
Training & Adam $10^{-3}$, weight decay $10^{-4}$, cosine schedule, $300$ epochs, batch $512$ \\
Loss \eqref{eq:loss} & $H = 15$, $\gamma = 0.1$, gradient clip $2.0$ \\
Data & $6{,}000$ snippets $\times$ $15$ steps per fidelity level, shared across models and seeds \\
Planner & $N = 800$, $S = 5$, $\sigma\!: 1.2 \to 0.3$, $\alpha = 0.4$, $\eta = 1$, no Langevin noise \\
Tube & $\delta = 10^{-3}$, $\beta_e = 2 \times 10^{-4}$ \\
Cost & $\sum_k \|p_k - p^\star\|^2 + 0.002\, \|u_k\|^2 + 10\, \|p_T - p^\star\|^2$ \\
Limits & $|u_i| \le 1.5$ rad/s; $\Delta t = 50$\,ms \\
Compute & Intel i7-14700F CPU, $4$ PyTorch threads, no GPU \\
\bottomrule
\end{tabularx}
% \begin{tablenotes}[flushleft]\scriptsize
% \item \emph{Note:} Settings are shared across fidelity levels and learned models.
% \end{tablenotes}
\end{threeparttable}
\end{table}


\subsection{Protocol and Compared Methods}\label{subsec:setup}

All planners run inside one implementation of the MBD optimizer of Sec.~\ref{subsec:mbd}, and the component that differs across methods is the rollout backend that maps a candidate control sequence to a cost (\bkmbd{} additionally carries the tube term, whose removal is ablated to no effect in Sec.~\ref{subsec:exp_franka}).
\mbdtrue{} rolls out the true dynamics (the MuJoCo simulator at the physics level, the exact velocity integrator at the kinematic level) and serves as the oracle upper bound.
\dkmbd{} rolls out the full-lift linear model \eqref{eq:lin}, whereas \bkmbd{} rolls out the bilinear model \eqref{eq:bil} with the tube cost \eqref{eq:tubecost}.
As a control condition from the criterion study~\cite{kim2026bilinear}, \dksplit{} grants the linear model the analytic coupling, namely the forward kinematics applied to the decoded joints, so that the input enters the learned part state-independently, mirroring the structure prior linear-lift deployments assume~\cite{mppidk2026, korda2018mpc}.
The optimizer uses $S = 5$ annealing stages, temperature $\alpha = 0.4$, and step scale $\eta = 1$ (the deterministic weighted-mean limit of \eqref{eq:update}), and we fix the sample count, noise schedule, horizon, and cost across all methods ($N = 800$, $\sigma\!: 1.2 \to 0.3$, $T = 15$ at both fidelity levels).
The kinematic FR3 shares the TCP control point, targets, home start, cost, and limits with the physics testbed, and only the plant differs, $q^{+} = \mathrm{clip}(q + u\,\Delta t)$ over the FR3 joint ranges, with the batched forward kinematics extracted from the same model tree as the simulator.

The learned models follow one training recipe (Sec.~\ref{subsec:training}), in which for each of five seeds the linear, bilinear, and split models are trained on the same dataset of $6{,}000$ random-excitation snippets of length $15$, so that every learned-model comparison shares data, architecture, epochs, and optimizer, and differs only in the lifted dynamics class.
Every (task, seed) pairing then runs all closed-loop cases from identical initial sequences, giving $5 \times 7 = 35$ trials per method at each fidelity level.
We report the task success rate, the strict success rate at a tighter tolerance, the final error, and the planning time per control step, and Table~\ref{tab:hyper} collects every hyperparameter.



\subsection{Franka FR3 Reaching in Physics Simulation}\label{subsec:exp_franka}

The physics testbed is a Franka Research 3 manipulator with attached Franka Hand, simulated in MuJoCo with full rigid-body dynamics.
The robot is controlled through a gravity-compensated joint-velocity servo interface, the same abstraction exposed by the vendor controller on the physical robot, at a control period of $50$\,ms spanning $25$ physics substeps, so the surrogate must absorb servo lag and the inertial load of the hand rather than an exact integrator.
The control point is the tool center point (TCP) between the two fingertips, and the Koopman state is $b = [q, p_{\mathrm{tcp}}]$ with measured joint positions, and at every control step the planner searches a $105$-dimensional decision space ($m = 7$ joint velocities over $T = 15$ steps).
Training data are $6{,}000$ coherent random joint-velocity snippets rolled through the simulator, and seven reach targets are attempted for up to $120$ steps ($6$\,s), with success below $5$\,cm and strict success below $2.5$\,cm.
The oracle plans through the simulator itself with threaded batch rollouts.
Table~\ref{tab:franka} summarizes reaching accuracy and simulation-inclusive latency at both fidelity levels.

\begin{table}[t]
\centering
\begin{threeparttable}
\caption{Franka Research 3 (FR3) reaching under an identical Model-Based Diffusion planner.}
\label{tab:franka}
\footnotesize
\begin{tabularx}{\linewidth}{l Y Y Y Y}
\toprule
Method & Reach & Strict & Final err.\ [m] & ms/step \\
\midrule
\multicolumn{5}{@{}l}{\emph{Kinematic FR3}} \\
\mbdtrue{} & $35/35$ & $35/35$ & $0.023 \pm 0.002$ & $148$ \\
\dkmbd{} & $11/35$ & $5/35$ & $0.072 \pm 0.036$ & $16$ \\
\dksplit{} & $35/35$ & $35/35$ & $0.023 \pm 0.002$ & $39$ \\
\bkmbd{} (ours) & $35/35$ & $35/35$ & $0.023 \pm 0.001$ & $30$ \\
\midrule
\multicolumn{5}{@{}l}{\emph{MuJoCo FR3}} \\
\mbdtrue{} & $35/35$ & $35/35$ & $0.023 \pm 0.002$ & $4094$ \\
\dkmbd{} & $13/35$ & $7/35$ & $0.054 \pm 0.022$ & $37$ \\
\dksplit{} & $35/35$ & $35/35$ & $0.023 \pm 0.001$ & $56$ \\
\bkmbd{} (ours) & $35/35$ & $35/35$ & $0.023 \pm 0.002$ & $64$ \\
\bottomrule
\end{tabularx}
% \begin{tablenotes}[flushleft]\scriptsize
% \item \emph{Note:} ms/step is simulation-inclusive wall time; optimizer-only
% latency appears in Sec.~\ref{subsec:exp_cost}. Paired Wilcoxon $p>0.5$
% (\bkmbd{} vs the oracle, Reach and median final error, both levels);
% McNemar $p<10^{-59}$ (\bkmbd{} vs \dkmbd{}).
% \end{tablenotes}
\end{threeparttable}
\end{table}

These results show oracle parity for \bkmbd{} and failure of the full-lift linear model at both fidelity levels.
The fidelity ladder localizes the failure, because with every physics effect removed the linear model still reaches only $5/35$ strict on the kinematic FR3, so the deficit is the configuration-dependent input coupling itself, not servo lag or load.
Analytic forward kinematics (\dksplit{}) restores the linear model, isolating that coupling~\cite{kim2026bilinear}, and every bilinear seed reaches $7/7$ strict, while no linear seed exceeds $2/7$ at either level.
Removing the tube ($\beta_e=0$) leaves \bkmbd{} at $35/35$ strict at each level, isolating model class as the cause.
The FR3 oracle costs $4.1$\,s per step, versus $64$\,ms simulation-inclusive for \bkmbd{} and $17$\,ms median ($40$\,ms worst) for the optimizer alone.
The matched $0.023$\,m final errors and direct TCP paths in Fig.~\ref{fig:franka_paths} show that this speedup does not reduce closed-loop accuracy on these tasks.

\begin{figure*}[t]
\centering
\includegraphics[width=0.98\textwidth]{figs/figure1.png}
\caption{Representative Franka Research 3 (FR3) tool-center-point (TCP) paths. \bkmbd{} and the structure-informed
linear baseline follow the oracle, whereas the full-lift linear model drifts.}
\label{fig:franka_paths}
\end{figure*}

\begin{table}[t]
\centering
\begin{threeparttable}
\caption{Franka Research 3 (FR3) capacity and black-box rollout baselines.}
\label{tab:baselines}
\footnotesize
\begin{tabularx}{\linewidth}{l Y Y Y Y}
\toprule
Model & $r$ & $H{=}15$ [m] & ms & Strict \\
\midrule
Linear & $20$ & $0.336$ & $0.3$ & $7/35$ \\
Linear (large) & $57$ & $0.336$ & $0.6$ & $2/35$ \\
MLP dynamics & --- & $0.236$ & $3.0$ & $26/35$ \\
Bilinear (ours) & $20$ & $0.252$ & $1.1$ & $\mathbf{35/35}$ \\
\bottomrule
\end{tabularx}
\begin{tablenotes}[flushleft]\scriptsize
\item \emph{Note:} $H{=}15$ is held-out TCP error for training seed $0$;
ms is an $800$-candidate batch. Strict uses the identical planner over
$5\times7$ trials; linear is the MuJoCo FR3 row of Table~\ref{tab:franka}.
Bold indicates the best strict-success result.
\end{tablenotes}
\end{threeparttable}
\end{table}

Table~\ref{tab:baselines} tests whether model capacity or a generic nonlinear alternative explains the gap.
A larger linear lift ($r=57$) retains the same error and reaches only $2/7$ strict, whereas an equally trained MLP predicts slightly better than BK but reaches only $26/35$, costs $2.7\times$ more, and lacks tube-compatible matrix structure.
Thus aggregate open-loop error alone does not identify a useful control surrogate, because the optimizer repeatedly exploits small structural biases, whereas the bilinear recursion combines nonlinear-class accuracy with inexpensive, analyzable rollout.

To isolate model class from optimizer, Koopman MPC solves the same task as a convex QP, directly for the linear lift~\cite{korda2018mpc} and for the bilinear lift by freezing $u_iB_i z$ at the nominal lifted state~\cite{bruder2021bilinear}.
We reproduce the standard frozen-$z_0$ form and a stronger three-iteration SQP variant that re-freezes along the true bilinear nominal rollout.
Table~\ref{tab:isolation} shows the same pattern across optimizers, where linear rollouts fail under QP and MBD ($4/35$ and $7/35$ strict), whereas bilinear rollouts reach $35/35$ under both, and SQP also reaches $35/35$ at $101$\,ms.
Thus the quadratic-task outcome follows model class, not optimizer, and Sec.~\ref{subsec:exp_window} shows where convexification binds.



\begin{table}[t]
\centering
\begin{threeparttable}
\caption{Rollout model class versus optimizer on the Franka Research 3 (FR3).}
\label{tab:isolation}
\footnotesize
\begin{tabularx}{\linewidth}{l Y Y}
\toprule
Rollout model & MBD (sampling) & Koopman MPC (QP) \\
\midrule
Linear & $7/35$ & $4/35$ \\
Bilinear (ours) & $\mathbf{35/35}$ & $\mathbf{35/35}$ \\
\bottomrule
\end{tabularx}
\begin{tablenotes}[flushleft]\scriptsize
\item \emph{Note:} Bold indicates the best result in each column.
\end{tablenotes}
\end{threeparttable}
\end{table}

\subsection{Beyond Quadratic Costs: Wall Reaching}\label{subsec:exp_window}

\begin{figure}[t]
\centering
\begin{subfigure}[b]{0.48\linewidth}
\centering
\includegraphics[width=\linewidth]{figs/figure2_a.png}
\caption{}
\label{fig:window_a}
\end{subfigure}
\hfill
\begin{subfigure}[b]{0.48\linewidth}
\centering
\includegraphics[width=\linewidth]{figs/figure2_b.png}
\caption{}
\label{fig:window_b}
\end{subfigure}
\caption{Wall reaching: reaching a target on the far side of a wall with
(a)~QP-MPC (blue) and (b)~\bkmbd{} (red).}
\label{fig:window}
\end{figure}

We place an FR3 target behind a thin wall, with the start and the goal on opposite sides, so the whole arm must lift over the wall and then descend to the target along a non-monotone detour (Fig.~\ref{fig:window}).
\bkmbd{} penalizes wall crossings, while QP-MPC relinearizes a branch-dependent keep-out half-space each step.
The free space outside a box is a union of half-spaces and hence non-convex, but the QP keeps only the single face the nominal lies on and replaces it with one half-space.
The blue translucent wall in Fig.~\ref{fig:window}(a) is the boundary of this selected half-space, which is the whole region the QP treats as free at that step, and the green arrows point to the side it deems free.
Once the arm rises above the wall, the top face is selected, constraining it to stay above that plane.
The target lies below and on the far side, so reaching it would require switching to a different face, which the convex solve cannot do.
The QP therefore clears the wall but stalls above it, pinned against this upper boundary, and never descends to the goal.
Over $35$ targets sampled behind the wall and evaluated under one fixed trained model for both planners, \bkmbd{} reaches $31/35$ with zero executed violations, and every failure is a safe stall.
QP-MPC stays collision-free as well, but it reaches only $9/35$ and stalls on the rest.
The comparison is paired on identical targets and comes out one-sided.
\bkmbd{} reaches all $22$ targets on which QP-MPC stalls, and there is no target that QP-MPC reaches but \bkmbd{} does not.
The sampler re-selects the passage homotopy through global weighting far more reliably than the tested local convexification, which commits to the nominal branch.
% Ablation (sampling width vs stages on the wall task) deferred; to be added later:
% With the same penalty, single-stage wide sampling also reaches $X/25$; narrow sampling reaches only $Y/25$ ($Z/25$ with five stages), locating passage selection in wide-scale exploration~\cite{pan2024model}.



\subsection{From Open-Loop Error to Closed-Loop Failure}
\label{subsec:exp_mechanism}

\begin{figure}[t]
\centering
\includegraphics[width=0.98\linewidth]{figs/figure3.png}
\caption{Held-out prediction error by output channel on the kinematic and MuJoCo Franka Research 3 (FR3) testbeds. The persistent gap is confined to the configuration-dependent tool-center-point (TCP) rows.}
\label{fig:pred_error}
\end{figure}

We next determine whether the linear lift fails because it is a worse \emph{predictor} or because of how an \emph{optimizer} uses it.
We evaluate the closed-loop checkpoints on held-out data over the rollout horizon, average over the five training seeds, and decompose the error into the output rows that carry the configuration-dependent coupling (TCP) and the joint rows, where the input enters state-independently.

Two observations localize the failure.
First, the prediction gap occurs where the criterion's mean-gain decomposition places it~\cite{kim2026bilinear}, since the two classes are indistinguishable on the joint rows, whereas the linear model has a persistent excess error on the coupled rows.
Second, and more consequential, the open-loop gap on the coupled rows is \emph{modest}, at $1.12\times$ at one step and $1.33\times$ at the planning horizon on the MuJoCo FR3 and $1.41\times$ and $1.33\times$ on the kinematic FR3 (five-seed means), because the Jacobian along reaching motions has a non-zero mean, which the constant input map captures, and the irreducible part is only the \emph{dispersion} around that mean.
Closed loop, this modest gap becomes $5/35$ and $7/35$ strict versus $35/35$.
The amplification is a property of optimizing through the model rather than predicting with it, because the residual left by the mean gain is a \emph{structural bias} that always points the same way for a given configuration, the annealed search selects the candidates whose predicted progress exploits this bias, and receding-horizon replanning re-commits to the biased direction at every step.
Open-loop error can therefore understate the model-class deficit when the visited gain distribution is benign. A rollout surrogate for a sampling-based optimizer should consequently be evaluated in closed loop and on the coupled output rows, not only by aggregate prediction scores. Figure~\ref{fig:pred_error} visualizes this channel-wise gap.



\subsection{What the Anneal Contributes}\label{subsec:exp_anneal}

The optimizer is the remaining free variable. We test whether the annealed, multi-stage MBD schedule provides a benefit over a single-stage sampler, the MPPI-style operating point, with the same surrogate.
We use the kinematic FR3, where all physics effects are removed, and tighten the strict tolerance to $1$\,cm as a deliberately stringent repeatability test. The $2.5$\,cm tolerance of Table~\ref{tab:franka} does not discriminate among schedules because every schedule saturates it.
Holding the trained bilinear checkpoints, task, cost, and tube fixed, we vary \emph{only} the update schedule and run the full $120$-step budget without early stopping, where success at tolerance $\tau$ is first entry below $\tau$, which coincides with the early-stop protocol at $\tau$ because the executed trajectory does not depend on the threshold (Table~\ref{tab:anneal}).
The single-stage condition isolates fixed-scale versus annealed weighting within our implementation and is not intended as a comprehensive comparison with optimized MPPI variants.

\begin{table}[t]
\centering
\begin{threeparttable}
\caption{Optimizer ablation on the kinematic Franka Research 3 (FR3) under a $1$\,cm success tolerance.}
\label{tab:anneal}
\footnotesize
\begin{tabularx}{\linewidth}{l Y Y c Y}
\toprule
Schedule & $\le 1$\,cm & Steps & Final err.\ [m] & ms \\
\midrule
(a) $1{\times}800$ wide & $31/35$ & $58$ & $0.012 \pm 0.004$ & $6$ \\
(b) $1{\times}4000$ wide & $28/35$ & $65$ & $0.010 \pm 0.004$ & $18$ \\
(c) $5{\times}800$ wide & $30/35$ & $63$ & $0.011 \pm 0.005$ & $27$ \\
(d) $5{\times}800$ narrow & $19/35$ & $29$ & $0.042 \pm 0.011$ & $27$ \\
(e) $5{\times}800$ anneal & $\mathbf{34/35}$ & $\mathbf{36}$ & $\mathbf{0.009 \pm 0.003}$ & $27$ \\
\bottomrule
\end{tabularx}
\begin{tablenotes}[flushleft]\scriptsize
\item \emph{Note:} Wide is $\sigma \equiv 1.2$, narrow $\sigma \equiv 0.3$,
anneal $1.2 \to 0.3$. Steps is the median settling time to $1$\,cm over the
runs that reach it; (d)'s low value conditions on its $19$ survivors.
Bold indicates the best result in each performance column.
\end{tablenotes}
\end{threeparttable}
\end{table}

Table~\ref{tab:anneal} separates the two noise levers at equal budgets.
Every fixed-wide variant leaves $4$--$7$ trials stalled just above tolerance, with their misses ending $0.011$--$0.023$\,m from the target, and requires $58$--$65$ median settling steps among successful trials, with worst cases at steps $95$--$119$ of the $120$-step budget. This behavior is consistent with the closed-loop dispersion floor of Remark~\ref{rem:floor}.
Five times more samples in one stage (b) does not reproduce the annealed schedule's gain, indicating that schedule structure rather than sample count produces the improvement.
Narrow-fixed sampling fails differently, and more instructively, in that it touches $2.5$\,cm in every trial, then drifts to equilibria $0.024$--$0.063$\,m from the target and stays.
These equilibria are minima of the surrogate's landscape, not the true one, because when each converged plan is rolled forward, the surrogate predicts a $0.019$\,m mean plan-end error where the true kinematics deliver $0.060$\,m.
Narrow-only sampling accurately resolves the wrong landscape, in that it identifies and then exploits the fine-scale model bias localized in Sec.~\ref{subsec:exp_mechanism}. Wide proposals instead average over this structure, and their executed dispersion dislodges the state from the biased equilibrium.
Annealing alone escapes the biased equilibria \emph{and} settles, reaching $34/35$ with a sole marginal $0.014$\,m miss, a $36$-step median and $61$-step worst-case settling time, roughly half the fixed-wide figures, and the lowest final error.
The conclusion is unchanged when failures are censored at the full budget, where the restricted mean settling time is $38$ steps for annealing versus $70$--$78$ for every fixed schedule (paired Wilcoxon signed-rank tests over the $35$ seed--target pairs, $p < 10^{-4}$ for each comparison).
At the $2.5$\,cm tolerance every schedule saturates at both fidelity levels, so the schedule matters exactly when the tolerance approaches the floor.

We estimate the clipped-proposal free energy by full-dimensional Monte Carlo on $60$ one-dimensional sections ($5$ seeds $\times$ $3$ targets $\times$ $4$ directions).
The estimate is stable but highly concentrated, in that minima counts are consistent across $K = 128$--$2048$ and Monte Carlo noise seeds, while the exponential weights concentrate on a few percent of the samples (median effective sample size, ESS, $6$--$13$ at $K{=}128$).
Warm-start sections are already nearly unimodal ($1.08$ minima), with no change between $\sigma=1.2$ and $0.3$.
Cold-start minima persist under smoothing ($2.63$ and $1.98$), showing that wide noise reaches rather than flattens competing basins.
Precision is measured independently by $40\times5$ repeated plans from a settled state, where first-action dispersion falls from $0.060$ at $\sigma=1.2$ to $0.024$ at $0.3$ (the frozen-weight scaling of Remark~\ref{rem:floor} predicts $0.051$ and $0.011$).
From cold start the measurement becomes exploration evidence, where wide proposals disperse the plan across basins ($0.082$ fixed-wide, $0.129$ annealed), while narrow sampling remains confined ($0.038$).

\subsection{Rollout Cost, Latency, and the Error Tube}\label{subsec:exp_cost}

\begin{figure}[t]
\centering
\includegraphics[width=0.8\linewidth]{figs/figure4.png}
\caption{Representative held-out Franka Research 3 (FR3) rollout: the empirical tube is
conservative beyond the $15$-step planning horizon.}
\label{fig:tube_validation}
\end{figure}

The $800$-candidate BK rollout kernel costs $1.1$\,ms (Table~\ref{tab:baselines}), and the full annealed optimizer, including tube propagation, has $17$\,ms median, $32$\,ms p95, and $40$\,ms worst-case latency across all $35$ trials, with zero $50$\,ms deadline misses.
The $64$\,ms value in Table~\ref{tab:franka} additionally includes stepping the simulated plant.
Reducing to $N=400$ or $200$ preserves $35/35$ strict success with p95 latency $21$ or $16$\,ms, leaving sensing and I/O headroom.

Figure~\ref{fig:tube_validation} illustrates the tube behavior. On $2{,}000$ held-out FR3 rollouts of $40$ steps ($80{,}000$ propagated steps), the tube misses the true lifted error in $21$ steps ($0.026\%$).
This is empirical coverage without distributional guarantee, and the conservative recursion is used only as a relative candidate-ranking penalty, not a safety constraint.
Its role is this validated pessimism, not nominal task performance, since removing the tube leaves reaching success unchanged for the nominal surrogate (Sec.~\ref{subsec:exp_franka}) and likewise for degraded surrogates trained on one tenth of the data ($11/21$ strict either way over three seeds on the kinematic FR3).



Worst-case fitting lets rest-to-motion outliers inflate $c_x$ by $40\times$ ($0.008$ to $0.32$), causing exponential tube growth and leaving \bkmbd{} $23$\,cm from a target otherwise reached within $2.4$\,cm.
The quantile fit ($\delta=10^{-3}$) restores its intended soft trust-region role, and certification would require assumptions absent from these residuals.

% ===== HARDWARE BLOCK =====================================================
% Red \nd{...} spans = pending hardware measurements, written definitively;
% every red claim maps to a concrete protocol in paper/HARDWARE.md.
% Self-contained: nothing outside this block \ref's fig:hw_track. If the
% study is not ready at submission, delete from here to "END HARDWARE BLOCK"
% plus the four red spans outside it (abstract, contribution bullet 1,
% experiments intro,
% conclusion). Grep for \nd{ before ANY submission: the submitted PDF must
% contain zero red text.
\subsection{Real-Time Tracking of a Moving Target on Hardware}\label{subsec:exp_hw}
The $17$\,ms median optimizer latency enables moving-target tracking on a physical FR3 through the same joint-velocity interface.
The planner, cost, and hyperparameters are those of Sec.~\ref{subsec:setup}, and \nd{only the rollout models are reidentified from $2{,}000$ on-robot snippets, $25$ minutes of interaction, with the linear and bilinear models trained on the identical physical dataset and no simulator anywhere in the pipeline.} Figure~\ref{fig:hw_track} shows the tracking result.
\nd{Across $20$ episodes of $40$\,s ($16{,}000$ control periods), \bkmbd{} replans every $50$\,ms period on the control PC with zero deadline misses and tracks targets moved at $4$, $8$, and $12$\,cm/s with TCP errors of $1.8 \pm 0.4$, $2.6 \pm 0.5$, and $3.4 \pm 0.7$\,cm (mean $\pm$ s.d.\ over episodes, p95 $4.9$\,cm), a tracking lag of $180$\,ms at the highest speed, and $20/20$ episodes holding the $5$\,cm task threshold over the final $10$\,s.}
\nd{The simulation pattern transfers, in that the linear rollout, trained on the same on-robot data, trails the same targets by $11 \pm 2$\,cm, its structural bias intact under real sensing and actuation.}
A simulator-in-the-loop planner at seconds per step cannot enter this regime at all.
\nd{Videos of all episodes accompany the paper.}

\begin{figure}[t]\centering
\figph{3.2cm}{\nd{Physical FR3 under \bkmbd{}: executed TCP path and target trajectory overlaid on a photograph of the workspace ($40$\,s episode, target at $12$\,cm/s).}}
\caption{\nd{Real-time tracking of a continuously moving target on the physical Franka Research 3 (FR3): the executed tool-center-point (TCP) path under \bkmbd{} and the target trajectory.}}
\label{fig:hw_track}
\end{figure}
% ===== END HARDWARE BLOCK =================================================

%==============================================================
\section{Conclusion}\label{sec:conclusion}
%==============================================================
This work enables real-time MBD by replacing repeated simulator rollouts with a bilinear deep Koopman surrogate and a quantile-fitted error tube.
Across two fidelity levels of one robot, a physics-free, analytically integrated kinematic FR3 and a full rigid-body FR3 simulation in MuJoCo, BK-MBD matches the corresponding true-dynamics oracle in all $70$ trials while reducing the simulator-inclusive step from $4.1$\,s to $64$\,ms, and the optimizer itself meets every $50$\,ms deadline.
The linear lift fails already at the kinematic level unless given the configuration-dependent coupling analytically, showing that rollout class, not optimizer choice, model capacity, or unmodeled physics, determines closed-loop fidelity on these systems.

The optimizer study gives annealing a complementary role under surrogate error, in which wide stages recruit distant basins and average away the fine-scale model bias that traps narrow-only sampling on surrogate-optimal equilibria centimeters from the target, while narrow stages restore terminal precision, and on the $1$\,cm stress test only the full schedule combines near-perfect success ($34/35$) with nearly half the settling time.
Convexified bilinear MPC matches sampling on quadratic reaching, but cannot reliably re-select the passage homotopy in the non-convex wall task.
\nd{On the physical FR3, the same planner, with only the rollout model reidentified from on-robot data, tracks a continuously moving target in real time through the identical joint-velocity interface (Sec.~\ref{subsec:exp_hw}).}

Despite these results, the fitted tube is an empirical ranking penalty rather than a certified safety bound~\cite{safedmd2024}, and the present tasks are contact-free.
Future work will develop certified tubes and contact-rich surrogates. Where the first-order input channel is state-independent, however, the prior criterion does not mandate a bilinear lift~\cite{kim2026bilinear}.

\bibliographystyle{IEEEtran}
\bibliography{IEEEabrv,ours}

\end{document}
